\documentclass{article}
\usepackage{graphicx}
\usepackage{amsmath}
\usepackage{booktabs}
\usepackage{siunitx}
\usepackage{float}

\title{Assignment 1 - EE4725 Quasi Optical Systems}
\author{Daniel Tyukov \\ Student Number: 5714699 \\ Microelectronics, TU Delft}
\date{February 2026}

\begin{document}

\maketitle

\tableofcontents
\newpage

%% ====================================================================
\section{Theory: Far-Field Evaluation Using the Spectral Green's Function}
%% ====================================================================

The electric field radiated by an electric current distribution can be written as a spectral-domain convolution:
\begin{equation}\label{eq:efield}
\vec{e}(x,y,z) = \frac{1}{4\pi^2}\int_{-\infty}^{\infty}\int_{-\infty}^{\infty}
\bar{\bar{G}}^{ej}(k_x,k_y)\,\vec{\tilde{J}}(k_x,k_y)\,
e^{-jk_x x}\,e^{-jk_y y}\,e^{-jk_z|z-z'|}\,dk_x\,dk_y
\end{equation}
where $\bar{\bar{G}}^{ej}$ is the dyadic spectral Green's function (SGF) for electric currents and $\vec{\tilde{J}}$ is the Fourier transform of the current distribution.

By evaluating the integral at the stationary phase point, the far field is obtained as:
\begin{equation}\label{eq:farfield}
\vec{E}^{far}(r,\theta,\phi) = j k_{zs}\,\bar{\bar{G}}^{ej}(k_{xs},k_{ys},z,z')\,
\vec{\tilde{J}}(k_{xs},k_{ys})\,\frac{e^{-jkr}}{2\pi r}
\end{equation}
where the stationary phase point corresponds to:
\begin{equation}
k_{xs} = k\sin\theta\cos\phi, \qquad
k_{ys} = k\sin\theta\sin\phi, \qquad
k_{zs} = k\cos\theta
\end{equation}

The dyadic SGF for electric currents in free space is:
\begin{equation}\label{eq:sgf}
\bar{\bar{G}}^{ej}(k_x,k_y) = -\frac{\zeta}{2kk_z}
\begin{bmatrix}
k^2-k_x^2 & -k_xk_y & -k_x(\pm k_z)\\
-k_xk_y & k^2-k_y^2 & -k_y(\pm k_z)\\
-k_x(\pm k_z) & -k_y(\pm k_z) & k^2-k_z^2
\end{bmatrix}
\end{equation}
where $+k_z$ is used for $z > z'$ and $-k_z$ for $z < z'$. The spectral variable $k_z$ is defined uniquely as:
\begin{equation}\label{eq:kz}
k_z = -j\sqrt{-(k^2-k_x^2-k_y^2)}
\end{equation}

The directivity is defined as:
\begin{equation}\label{eq:directivity}
D(\theta,\phi) = \frac{U(\theta,\phi)}{P_{\text{rad}}/(4\pi)}
\end{equation}
where the radiation intensity is:
\begin{equation}\label{eq:radintensity}
U(\theta,\phi) = \frac{|\vec{E}^{far}(r,\theta,\phi)|^2}{2\zeta}\,r^2
\end{equation}
and the total radiated power is:
\begin{equation}\label{eq:prad}
P_{\text{rad}} = \int_0^{2\pi}\int_0^{\pi} U(\theta,\phi)\sin\theta\,d\theta\,d\phi
\end{equation}

%% ====================================================================
\section{MATLAB Implementation}
%% ====================================================================

Four reusable MATLAB functions were implemented:

\begin{enumerate}
\item \texttt{EJ\_SGF(er, k, kx, ky)} -- computes the nine components of the spectral dyadic Green's function $\bar{\bar{G}}^{ej}$.
\item \texttt{FTCurrent(k0, er, kx, ky, l, w)} -- computes the analytical Fourier transform of the dipole current distribution.
\item \texttt{farfield(k0, R\_FF, TH, PH, kz, Gxx, Gyx, Gzx, Jx)} -- evaluates the far field via the stationary phase formula~\eqref{eq:farfield}. The product $jk_{zs}\bar{\bar{G}}^{ej}$ is formed by multiplying the passed SGF components with $k_z$; at $\theta=\pi/2$ where $k_z\to 0$ and the SGF diverges, the resulting $0\times\infty$ is replaced by the known analytical limit $-j\zeta/(2k)$ times the corresponding matrix element.
\item \texttt{Directivity(E\_tot, Theta, dth, dph, er, r)} -- computes the directivity and total radiated power by numerical integration.
\end{enumerate}

%% ====================================================================
\section{Q1: Elementary Electric Source (3 points)}
%% ====================================================================

\subsection{Problem Setup}

An $\hat{x}$-oriented elementary electric source is placed at the origin of the coordinate system:
\begin{equation}
\vec{j}(x,y,z) = \delta(x)\,\delta(y)\,\delta(z)\,\hat{x}
\end{equation}
Its Fourier transform is simply $\vec{\tilde{J}}(k_x,k_y)=\hat{x}$.

Since the current is $x$-directed, only the first column of the dyadic SGF is relevant:
\begin{equation}
\vec{E} = \bar{\bar{G}}^{ej}\cdot\hat{x} =
\begin{bmatrix} G_{xx} \\ G_{yx} \\ G_{zx} \end{bmatrix}
\end{equation}

The parameters used are $f=\SI{30}{GHz}$, $k_x\in[0,\,3k_0]$, and $k_y=0$.

\subsection{Q1.1: SGF Components}

Figure~\ref{fig:q1_sgf} shows the real and imaginary parts of the three SGF components $G_{xx}$, $G_{yx}$, and $G_{zx}$ as a function of $k_x/k_0$.

\begin{figure}[H]
\centering
\includegraphics[width=0.95\textwidth]{Q1_SGF_components.png}
\caption{Real and imaginary parts of the SGF components $G_{xx}$, $G_{yx}$, and $G_{zx}$ for an $\hat{x}$-oriented elementary source at $f=\SI{30}{GHz}$, $k_y=0$. The dashed vertical line marks $k_x=k_0$.}
\label{fig:q1_sgf}
\end{figure}

Key observations from the plots:

\begin{itemize}
\item \textbf{$G_{xx}$}: real and negative for $k_x<k_0$, purely imaginary and positive for $k_x>k_0$. The singularity at $k_x=k_0$ corresponds to $k_z=0$.
\item \textbf{$G_{yx}$}: identically zero for all $k_x$ because $k_y=0$, so the off-diagonal term $-k_xk_y=0$.
\item \textbf{$G_{zx}$}: purely real over the entire range. This is because the $k_z$ in the numerator ($-k_xk_z$) cancels the $k_z$ in the denominator of the prefactor, giving $G_{zx}=\zeta k_x/(2k)$, which is real regardless of whether $k_z$ is real or imaginary.
\end{itemize}

\subsection{Q1.2: Imaginary Region of the $x$-Component}

The $x$-component of the SGF is:
\begin{equation}
G_{xx} = -\frac{\zeta}{2kk_z}\left(k^2-k_x^2\right)
\end{equation}

\begin{itemize}
\item \textbf{Propagating region} ($k_x < k_0$, i.e.\ $k_x^2+k_y^2<k^2$): $k_z=\sqrt{k^2-k_x^2}$ is real and positive. The prefactor $-\zeta/(2kk_z)$ is real, and $(k^2-k_x^2)$ is real and positive. Therefore $G_{xx}$ is \textbf{purely real} (and negative).

\item \textbf{Evanescent region} ($k_x > k_0$, i.e.\ $k_x^2+k_y^2>k^2$): $k_z=-j\sqrt{k_x^2-k^2}$ is purely imaginary. The prefactor becomes $-\zeta/(2k\cdot(-j|\cdot|))$, which is purely imaginary. Since $(k^2-k_x^2)$ remains real (now negative), the product $G_{xx}$ is \textbf{purely imaginary}.
\end{itemize}

This is confirmed by Figure~\ref{fig:q1_sgf}: the real part of $G_{xx}$ (blue solid line) is nonzero only for $k_x/k_0<1$, while the imaginary part (red dashed line) is nonzero only for $k_x/k_0>1$.

%% ====================================================================
\section{Q2: Far Field of a Dipole (4 points)}
%% ====================================================================

\subsection{Problem Setup}

An $\hat{x}$-directed dipole antenna is centered at the origin with length $L$ along $x$ and width $W$ along $y$, radiating in free space.

The spatial current distribution is modeled using the open-circuited transmission-line approximation:
\begin{equation}
\vec{J}_x(x,y) = l(x)\,t(y)\,\hat{x}
\end{equation}
with Fourier transform $\tilde{J}_x^{FT}(k_x,k_y)=L(k_x)\,T(k_y)$, where:

\begin{itemize}
\item Transverse distribution:
\begin{equation}
t(y)=\frac{1}{w}\,\text{rect}\!\left(\frac{y}{w}\right),\qquad
T(k_y)=\text{sinc}\!\left(\frac{k_y w}{2}\right)
\end{equation}

\item Longitudinal distribution:
\begin{equation}\label{eq:Lkx}
l(x)=\frac{\sin\!\bigl(k_{eq}(\tfrac{l}{2}-|x|)\bigr)}{\sin(k_{eq}\tfrac{l}{2})},\qquad
L(k_x)=\frac{2k_{eq}\bigl(\cos(\tfrac{k_x l}{2})-\cos(\tfrac{k_{eq}l}{2})\bigr)}{(k_{eq}^2-k_x^2)\sin(k_{eq}\tfrac{l}{2})}
\end{equation}
\end{itemize}

Per the professor's clarification, $k_{eq}=k_0$ (the free-space wavenumber).

\subsection{Q2.1: 1D Far-Field Plots in Main Planes}

The dipole dimensions are $L=\lambda_0/2=\SI{5}{mm}$ and $W=\lambda_0/40=\SI{0.25}{mm}$ at $f=\SI{30}{GHz}$ ($\lambda_0=\SI{10}{mm}$). The total far field $|\vec{E}^{far}|=\sqrt{|E_\theta|^2+|E_\phi|^2}$ is computed for $\phi=0^\circ$, $45^\circ$, and $90^\circ$.

\begin{figure}[H]
\centering
\includegraphics[width=0.85\textwidth]{Q2_1_farfield_1D.png}
\caption{Normalized far-field pattern of the half-wave dipole in the three main planes.}
\label{fig:q2_1d}
\end{figure}

Observations from Figure~\ref{fig:q2_1d}:
\begin{itemize}
\item \textbf{$\phi=0^\circ$ (E-plane, $xz$-plane)}: The pattern drops sharply towards $\theta=90^\circ$ (along the dipole axis), reaching a deep null. This is characteristic of a dipole: no radiation along its own axis.
\item \textbf{$\phi=90^\circ$ (H-plane, $yz$-plane)}: The pattern is nearly constant (omnidirectional in the H-plane), as expected for a thin dipole. The slight drop at large $\theta$ is due to the finite width $W$.
\item \textbf{$\phi=45^\circ$}: An intermediate pattern between the E- and H-planes.
\item All three planes share the same broadside ($\theta=0^\circ$) level, confirming the maximum radiation is in the $+z$ direction.
\end{itemize}

\subsection{Q2.2: UV Representation of the Far Field}

The far field is plotted in the UV plane ($U=\sin\theta\cos\phi$, $V=\sin\theta\sin\phi$) with a 10\,dB dynamic range.

\begin{figure}[H]
\centering
\includegraphics[width=0.7\textwidth]{Q2_2_farfield_UV.png}
\caption{Far-field pattern of the half-wave dipole in the UV plane (10\,dB dynamic range). The $U$-axis corresponds to the dipole axis direction.}
\label{fig:q2_uv}
\end{figure}

The UV representation in Figure~\ref{fig:q2_uv} clearly shows:
\begin{itemize}
\item The maximum radiation is at the center ($U=V=0$, broadside).
\item The pattern exhibits a null along the $U$-axis ($U=\pm1$, $V=0$), i.e.\ along the dipole axis ($\hat{x}$-direction at $\theta=90^\circ$, $\phi=0^\circ$ or $180^\circ$).
\item The radiation is nearly uniform along the $V$-axis, consistent with the omnidirectional H-plane pattern.
\item The circular boundary ($U^2+V^2=1$) corresponds to $\theta=90^\circ$ (the horizon).
\end{itemize}

\subsection{Q2.3: Broadside Directivity vs.\ Frequency}

For a fixed dipole with $L=\SI{5}{mm}$ and $W=\SI{0.25}{mm}$, the broadside directivity is computed as a function of frequency from \SI{5}{GHz} to \SI{60}{GHz}. The total radiated power is obtained by integrating over the upper hemisphere and multiplying by 2 (exploiting the symmetry of the free-space pattern about $z=0$).

\begin{figure}[H]
\centering
\includegraphics[width=0.85\textwidth]{Q2_3_directivity_vs_freq.png}
\caption{Broadside directivity of the dipole ($L=\SI{5}{mm}$, $W=\SI{0.25}{mm}$) as a function of frequency. The dashed red line marks the theoretical half-wave dipole directivity of \SI{2.15}{dBi}.}
\label{fig:q2_dir}
\end{figure}

Observations from Figure~\ref{fig:q2_dir}:
\begin{itemize}
\item At $f=\SI{30}{GHz}$ (where $L=\lambda_0/2$, i.e.\ half-wave resonance), the broadside directivity is $\approx\SI{2.12}{dBi}$, matching the theoretical value for a half-wave dipole (\SI{2.15}{dBi}).
\item At low frequencies ($f\ll\SI{30}{GHz}$), the dipole is electrically short ($L\ll\lambda$) and approaches the Hertzian dipole limit with $D\approx\SI{1.76}{dBi}$.
\item At high frequencies ($f\gg\SI{30}{GHz}$), the dipole becomes electrically long ($L>\lambda$), and the directivity increases as the pattern narrows.
\end{itemize}

%% ====================================================================
\section{Q3: Dipole with a Backing Reflector (3 points)}
%% ====================================================================

\subsection{Problem Setup}

The same dipole ($L=\lambda_0/2$, $W=\lambda_0/40$ at \SI{30}{GHz}) is now placed at height $h$ above an infinite PEC ground plane. By the image theorem for a perfect electric conductor ($\hat{n}\times\vec{E}=0$):

\begin{itemize}
\item The real source at $z=+h$ has current $+\hat{x}$.
\item The image source at $z=-h$ has current $-\hat{x}$ (tangential electric current reverses sign).
\end{itemize}

The total far field is computed using the free-space SGF applied to both sources. The combined current spectrum includes an array factor:
\begin{equation}\label{eq:arrayfactor}
\vec{\tilde{J}}_{\text{total}}(k_{xs},k_{ys}) = \vec{\tilde{J}}(k_{xs},k_{ys})\left[e^{jk_{zs}h}-e^{-jk_{zs}h}\right]
= \vec{\tilde{J}}(k_{xs},k_{ys})\cdot 2j\sin(k_{zs}\,h)
\end{equation}
where $k_{zs}=k_0\cos\theta$ at the stationary phase point.

\subsection{Q3.1: Far-Field Comparison with Free-Space Dipole}

The far field is computed for $h=\SI{7.5}{mm}=3\lambda_0/4$ and compared with the free-space dipole pattern from Q2.

\begin{figure}[H]
\centering
\includegraphics[width=0.95\textwidth]{Q3_1_farfield_comparison.png}
\caption{Far-field comparison between the free-space dipole and the dipole with PEC reflector at $h=\SI{7.5}{mm}$, for $\phi=0^\circ$, $45^\circ$, and $90^\circ$. Amplitudes are normalized to the free-space broadside level.}
\label{fig:q3_comparison}
\end{figure}

Observations from Figure~\ref{fig:q3_comparison}:
\begin{itemize}
\item The PEC reflector increases the broadside radiation by approximately \SI{6}{dB} relative to the free-space case, because the image source constructively interferes at broadside when $h=3\lambda_0/4$ (the array factor $2\sin(k_0 h\cos\theta)$ evaluated at $\theta=0$ gives $2\sin(3\pi/2)=-2$, so $|AF|=2$, doubling the field and adding \SI{6}{dB}).
\item The reflector introduces additional nulls in the pattern due to destructive interference between the real and image sources at specific angles.
\item At $\theta=90^\circ$ (the horizon), the array factor $2j\sin(k_0 h\cdot 0)=0$, producing a null for all phi-cuts. This is expected: the PEC ground plane enforces zero tangential electric field at grazing incidence.
\item In the H-plane ($\phi=90^\circ$), the free-space pattern is nearly flat, so the interference lobes from the reflector are clearly visible.
\end{itemize}

\subsection{Q3.3: Broadside Directivity vs.\ Height $h$}

The broadside directivity (linear scale) is plotted as a function of the distance $h$ from $\lambda_0/10$ to $2\lambda_0$. For the PEC case, the total radiated power is integrated over the upper hemisphere only ($\theta\in[0,\pi/2]$), since no radiation exists below the ground plane.

\begin{figure}[H]
\centering
\includegraphics[width=0.85\textwidth]{Q3_3_directivity_vs_h.png}
\caption{Broadside directivity (linear scale) as a function of the dipole height $h/\lambda_0$ above the PEC ground plane.}
\label{fig:q3_dir}
\end{figure}

Observations from Figure~\ref{fig:q3_dir}:
\begin{itemize}
\item The directivity oscillates with $h$, showing clear nulls at $h=n\lambda_0/2$ for $n=1,2,3,\ldots$ (marked by red dashed lines). Between the nulls, the directivity peaks at values around 6--7 (linear), corresponding to a significant enhancement over the free-space dipole ($D\approx1.64$).
\item At very small $h$ (near $\lambda_0/10$), the directivity is particularly high ($\approx7.5$) because the array factor $2j\sin(k_0 h\cos\theta)$ concentrates radiation near broadside when $h$ is small.
\end{itemize}

\subsubsection*{Integration domain for radiated power}

Since the image theorem replaces the PEC with image sources, the resulting fields are valid only in the upper half-space ($z>0$). Therefore, the total radiated power is computed by integrating over the upper hemisphere:
\begin{equation}
P_{\text{rad}} = \int_0^{2\pi}\int_0^{\pi/2} U(\theta,\phi)\sin\theta\,d\theta\,d\phi
\end{equation}

\subsubsection*{Null at broadside}

At broadside ($\theta=0$), $k_{zs}=k_0$. The array factor is $2j\sin(k_0 h)$, which vanishes when:
\begin{equation}
k_0 h = n\pi \quad\Longrightarrow\quad h = \frac{n\lambda_0}{2}, \qquad n=1,2,3,\ldots
\end{equation}
This is confirmed numerically in Figure~\ref{fig:q3_dir}, where the nulls appear precisely at $h/\lambda_0=0.5$, $1.0$, $1.5$, and $2.0$.

\subsubsection*{Power radiated for $h=0$}

When $h=0$, the image source at $z=-h=0$ coincides spatially with the real source. Since the image current is $-\hat{x}$ (opposite to the real current $+\hat{x}$), the two sources cancel exactly:
\begin{equation}
\text{AF}\big|_{h=0} = 2j\sin(k_0\cdot 0) = 0
\end{equation}
Therefore the total field is zero everywhere and $P_{\text{rad}}=0$. This was verified numerically: $P_{\text{rad}}(h=0)=\num{0.0}\,\si{W}$.

\end{document}
